\paragraph{Scalarity dichotomy and noncentral condensate selection of the leaf bivector.} Let \(\hat X^I\) be self-adjoint coordinate operators generating a unital operator algebra, and call \(\Theta^{IJ}_\omega:=-\mathrm i\,\omega([\hat X^I,\hat X^J])\), indexed by normalized states \(\omega\), the state-induced bivector. This datum is independent of the state if and only if every commutator \([\hat X^I,\hat X^J]\) is scalar; in finite dimension the trace forces such a scalar to vanish, so the centrality mechanism of the constant-\(\Theta\) chart-blindness is exactly the scalar case of the dichotomy, and every finite-\(N\) matrix regularization with noncommuting coordinates realizes the complement: the commutator is traceless, hence non-scalar, and the state-induced bivector has nontrivial spread over states (block-local states in a direct sum of \(N=5\) and \(N=8\) 't Hooft blocks give distinct bivector values). In the Lie-type class \([\hat X^I,\hat X^J]=\mathrm i\,c^{IJ}{}_K\hat X^K\), possibly with state-independent scalar terms, the bivector of a state with one-point functions \(x^K=\omega(\hat X^K)\) equals \(c^{IJ}{}_Kx^K\) up to the constant contributed by the scalar terms, a derived leaf-dependent bivector field: for the spin-\(j\) fuzzy sphere \(\hat X^i=\ell\hat J^i\) with \([\hat X^1,\hat X^2]=\mathrm i\ell\hat X^3\), the bivector \(\Theta^{12}_\omega=\ell\,\omega(\hat X^3)\) takes values spanning \([-\ell^2 j,\ell^2 j]\) over pure states, its coherent-state fluctuation has relative width exactly \(1/\sqrt j\) about the mean bivector, in the commutator-fluctuation-over-bivector-norm measure, and the exact Casimir identity \(\hat X\cdot\hat X=\ell^2 j(j+1)\mathbb 1\) fixes the leaf radius, the bivector field being the Kirillov--Kostant--Souriau field of that orbit. Where the derived bivector field is nondegenerate on the leaf, the set of chart 1-jets carrying it to the constant ambient bivector is a single orbit of the symplectic group (dimension \(10\) at leaf dimension \(4\)): the bivector field supplies the gauge-invariant content of the semiclassical-chart 1-jet, with the bivector-automorphism group (\(\mathfrak{so}(3)\), dimension \(3\), for the fuzzy sphere) as the residual gauge. Where it is degenerate (rank \(2\) at every point of the fuzzy sphere, hence at every point of any four-coordinate embedding of a two-dimensional orbit), rank invariance under frame change blocks this supply, and characteristic-direction data beyond the bivector are required. These are exact linear-algebraic and \(\mathfrak{su}(2)\)-representation identities with explicit finite-dimensional matrix computations; they do not compute the order-\(\Theta^2\) correction selected in the Lie-type class, they do not determine whether a derived bivector germ fixes the embedding 2-jet within the \(20\)-dimensional fiber of embedding 2-jets at fixed bivector 1-jet, they do not exhibit a noncentral condensate of the matrix model of this book, they do not cover generic non-Lie-type condensates, and this paragraph does not close Problem~A.
